\section{Results}

\subsection{Part 1: Baseline on the RAG evaluation set}
On the 23 single-article questions the baseline used 2.0 tool calls, 5.9\,s and \$0.019 per question. It
retrieved the exact source article for 12 of the 19 linked questions (63\%). Manual inspection showed that
five of the seven misses were generic questions (``What is Agrivoltaics?'', ``What is the EU's Green Deal
Industrial Plan?'') that many articles answer. The agent answered them correctly from another article; for the
Green Deal question, the article it cited gives the reference definition almost word for word. The two real misses
needed one specific detail, such as the origin of Iceland's geothermal turbines. Here the agent stated that it could
not find the answer rather than inventing one. Article-level recall therefore understates answer quality on
generic questions, a caveat that also applies to the 50-question set.

\subsection{Retrieval accuracy}
Table~\ref{tab:retrieval} ranks articles for the 50 synthesis questions. The from-scratch BM25 is the
strongest single retriever (MRR 0.636), well ahead of the bge-base dense model (0.542). These questions are
dense with proper nouns and figures (\textit{Kosice}, \textit{CMBlu}, \textit{EUR 56.1 million}), which is
exactly what exact lexical match rewards. A general 110M-parameter embedding instead places ``Kosice
geothermal'' near every geothermal article in the corpus. Plain RRF does not beat BM25, because it gives the weaker
dense list an equal vote. The cross-encoder recovers that loss and gives the best Recall@5 (0.607), Recall@10
(0.673) and nDCG@10 (0.576) of the single-query retrievers.

\begin{table}[t]
\caption{Ranked article retrieval on CleanTech-50 (mean over 50 questions)}
\label{tab:retrieval}
\centering
\setlength{\tabcolsep}{4pt}
\begin{tabular}{lcccccc}
\toprule
Configuration & R@1 & R@3 & R@5 & R@10 & MRR & nDCG \\
\midrule
BM25 (scratch)          & \textbf{.355} & \textbf{.530} & .550 & .630 & \textbf{.636} & .560 \\
Dense (bge-base)        & .330 & .410 & .460 & .543 & .542 & .476 \\
Hybrid RRF              & .337 & .493 & .545 & .627 & .610 & .542 \\
Hybrid RRF + rerank     & .352 & .527 & \textbf{.607} & .673 & .619 & \textbf{.576} \\
\midrule
Crew v1: planner queries  & .317 & .382 & .517 & .637 & .557 & .512 \\
Crew v2: anchor + planner & .295 & .507 & .573 & \textbf{.697} & .595 & .556 \\
\bottomrule
\end{tabular}
\end{table}

The bottom rows rank the articles retrieved by the crew's own search queries, fused with RRF. v1's
decomposed queries put fewer referenced articles in the top three than searching with the plain question (R@3 0.38
vs 0.53). Turning ``barriers mentioned in \emph{the interviews}'' into generic keyword queries drops the words that
point to the specific interview. Adding an anchor query on the whole question (v2) restores the top ranks and gives
the best Recall@10 and Recall@20 (0.724) of any configuration, because the sub-queries reach articles for the second
half of comparison questions.

Fig.~\ref{fig:retrieval} (right) shows what the agents read and cite. The crew retrieves the most referenced
articles (evidence recall 0.766 vs 0.657 for the baseline) and cites the most (0.714 vs 0.645). Swapping only the
retriever inside the baseline closes about a third of that gap (0.698). The crew's weakness is breadth: it reads
about 34 articles and cites about 7 per answer, so its citation precision (0.290) is lower than the baseline's
(0.392).

\begin{figure*}[t]
\centering
\includegraphics[width=0.85\textwidth]{retrieval.png}
\caption{Left: Recall@$k$ of the ranked retrievers. Right: evidence recall, citation recall and citation precision of
the agents.}
\label{fig:retrieval}
\end{figure*}

\subsection{Response accuracy}
Table~\ref{tab:main} and Fig.~\ref{fig:judge} give the main results. The judge metrics carry 95\% bootstrap confidence intervals in the
notebook; with 50 questions, fact-recall differences below about 0.06 are within noise.

\begin{table*}[t]
\caption{Response accuracy on CleanTech-50 (mean over 50 questions). FR = judge fact recall; Corr = judge
correctness (1--5); Faith = share of claims supported by the cited sources; Contr = contradicted key facts;
BERTS = rescaled BERTScore F1; cost and latency per question.}
\label{tab:main}
\centering
\begin{tabular}{lcccccccc}
\toprule
System & FR $\uparrow$ & Corr $\uparrow$ & Faith $\uparrow$ & Contr $\downarrow$ & BERTS $\uparrow$ & ROUGE-L $\uparrow$ & Cost (\$) & Latency (s) \\
\midrule
Closed-book (Sonnet 5.5)       & .474 & 3.06 & --   & .010 & .060 & .139 & .007 & 6.1 \\
Baseline (Part 1, ReAct)       & .672 & 3.76 & \textbf{.966} & .016 & .111 & \textbf{.176} & .022 & 6.9 \\
Baseline + hybrid retrieval    & .628 & 3.76 & .964 & .013 & \textbf{.112} & .172 & .022 & 9.4 \\
Crew v1                        & .559 & 3.40 & .927 & .029 & .070 & .160 & .173 & 55.8 \\
Crew v2                        & .603 & 3.62 & .926 & .040 & .080 & .164 & .175 & 43.5 \\
\quad v2, no critic (first draft) & .630 & 3.66 & .909 & .031 & .094 & .170 & -- & -- \\
\quad v2, plain prose (no STE) & \textbf{.683} & \textbf{3.84} & .915 & .024 & .097 & .165 & -- & -- \\
\bottomrule
\end{tabular}
\end{table*}

\begin{figure*}[t]
\centering
\includegraphics[width=0.92\textwidth]{judge_scores.png}
\caption{LLM-judge scores per system with 95\% bootstrap confidence intervals. Closed-book has no citations, so no faithfulness.}
\label{fig:judge}
\end{figure*}

\textbf{Retrieval is the largest effect.} Grounding the same model in the corpus raises fact recall from 0.474
to 0.672 and correctness from 3.06 to 3.76 (both $p<0.001$; 35 wins, 8 losses). The per-topic view
(Fig.~\ref{fig:topics}) shows why. Without retrieval, recall collapses on Geothermal (0.27), Policy (0.23) and
Solar (0.33), whose questions hinge on specific 2023 projects that no model saw in pre-training. It stays at 0.71 on
Climate \& Environment, where general science suffices.

\textbf{Answers are faithful.} 97\% of the baseline's claims and 93\% of the crew's are supported by the passages
they cite, and at most 4\% of key facts are contradicted. The output guardrail removed no hallucinated citation
in any answer: the writer cited only ids it had been shown.

\textbf{The multi-agent crew does not beat the single agent.} Crew v2 recalls fewer reference facts than the
baseline (0.603 vs 0.672, $p=0.002$) and is less faithful ($p<0.001$); the correctness gap is not significant
($p=0.07$). Its best configuration, the first draft re-written without the STE instruction, is the top system
on fact recall (0.683) and correctness (3.84), but it is only level with the baseline, at eight times the cost and
six times the latency.

\begin{figure}[t]
\centering
\includegraphics[width=\columnwidth]{topics.png}
\caption{Judge fact recall by topic and system (number of questions in parentheses).}
\label{fig:topics}
\end{figure}

\subsection{Ablations}
Because both ablations re-use v2's saved evidence, each changes exactly one component.

\textbf{ASD-STE100 trades content for readability.} The STE instruction works as intended
(Table~\ref{tab:ste}). The crew's sentences average 11 words against 23 in the references, 97\% fit the STE
limit, passive voice falls to 2\%, and Flesch Reading Ease rises from 16 (``very difficult'') to 50. On identical
evidence, however, the plain-prose writer recalls 5.2 points more key facts ($p<0.001$) and scores 0.18 higher on
correctness ($p=0.01$). Short sentences leave room for fewer facts in the same word budget, and STE's simplification
drops qualifiers that the reference counts as facts. This is a concrete, measured cost of Karpathy's suggestion:
STE answers are easier to \emph{review}, but on this task they are about 8\% less complete.

\begin{table}[t]
\caption{Readability and STE compliance (rule-based checker)}
\label{tab:ste}
\centering
\setlength{\tabcolsep}{4pt}
\begin{tabular}{lcccc}
\toprule
Text & Words/sent. & $\le$ limit & Passive & Flesch \\
\midrule
Reference answers      & 23.0 & 62\% & 10.0\% & 16 \\
Closed-book            & 16.8 & 83\% & 9.3\% & 31 \\
Baseline (Part 1)      & 12.1 & 99\% & 8.6\% & 43 \\
Crew v2 (STE)          & 11.0 & 97\% & 2.3\% & 50 \\
Crew v2, plain prose   & 16.8 & 84\% & 14.4\% & 37 \\
\bottomrule
\end{tabular}
\end{table}

\textbf{The critic adds cost, not measurable quality.} Comparing v2's final answer with its first draft, the critic
loop has no significant effect on fact recall ($-0.028$, $p=0.13$) or correctness ($p=0.56$). It nudges
faithfulness up ($+0.017$, $p=0.08$), which is its purpose. It requested a revision for 44 of 50 first drafts, and
most revisions added caveats rather than facts.

\textbf{Retrieval alone does not transfer to the baseline.} Giving the baseline the hybrid retriever raised its
evidence recall but left its answers unchanged (fact recall $-0.04$, $p=0.09$; correctness identical). A single
focused query per turn already finds what the agent needs for its answer; the extra articles mattered only in the
crew, where several researchers query different aspects.

\subsection{From v1 to v2: error analysis}
\label{sec:v2}
Crew v1 scored clearly below the baseline (fact recall 0.559). Inspecting its traces revealed three
implementation faults and two design faults:
(i) text read through \texttt{get\_article} or \texttt{summarize\_article} was cited but never recorded as
evidence;
(ii) the writer saw only the first 24 passages \emph{in arrival order}, so the last researcher's evidence could be
dropped;
(iii) the evaluation itself truncated cited sources to 9{,}000 characters, which penalised answers with many
citations (fixing it raised every system's measured faithfulness, the baseline's from 0.85 to 0.97);
(iv) decomposed queries lost the question's specificity (R@3 0.38); and (v) the critic demanded tangential
additions. v2 records all evidence, ranks passages by reranker score, adds the anchor query, and narrows the
critic. Fact recall rose to 0.603 ($p=0.08$) and correctness from 3.40 to 3.62 ($p=0.02$). Since these fixes were
motivated by errors on the test set itself, v2's scores are optimistic; v1 is the fair first-attempt result.

\subsection{Judge reliability and cost}
The judge's fact recall correlates with BERTScore at $\rho=0.54$ and with ROUGE-L at $\rho=0.48$: it tracks meaning,
but it is not simply a lexical metric. Re-grading 25 random answers with fresh calls gave the same correctness score
80\% of the time and never differed by more than one point (quadratic-weighted $\kappa=0.82$). All agent runs in this study
(Part~1, seven systems and both crew versions) cost \$24.67 in Anthropic API calls, as recorded by the usage
callbacks; the judge's OpenAI calls are billed separately.

\section{Discussion}
\textbf{Why does the crew not win?} The error analysis points to a mismatch between the system's search space and
the benchmark's construction. Every reference answer synthesizes 2--5 articles published in the same week, and
questions refer to them as ``the interviews and reports''. The crew searches 20{,}111 articles with several
researchers, so it finds the referenced articles more often than any other system. It then writes a broad,
well-supported survey from about seven sources, and the reference facts are crowded out. The three lowest-scoring
crew answers all follow this pattern. For Q21, the crew read the Diverso interview but answered ``barriers to
geothermal'' from thirteen permitting articles, and the judge rated it accurate yet off-target. The single ReAct
agent, by contrast, issues one or two focused queries and stays close to whatever it finds first, which on this
benchmark is usually the right article.

\textbf{What would help without overfitting?} A planner step that first resolves \emph{which} documents a question
refers to (by entity, outlet and date) and then restricts synthesis to them; a critic that runs only when the
draft contains unsupported claims; and STE applied as a presentation layer (a summary box) on top of complete
prose, instead of as a constraint on the answer itself.

\textbf{Limitations.} Fifty questions give wide confidence intervals; one judge model scores against one reference
per question, and its key facts come from the same judge; and article-level retrieval metrics penalise correct
answers drawn from other articles. Agent runs are stochastic (no temperature control on these models), so a
re-run would shift individual scores; the reported numbers come from one cached run.

\section{Conclusion}
I built an agentic RAG system for 20{,}111 cleantech articles with six extensions: a from-scratch BM25 with RRF and
reranking, an MCP tool server with a hand-written client, a planner/researcher/writer/critic crew with memory,
input and output guardrails, ASD-STE100 answers with a compliance checker, and scholarly search. A rigorous
evaluation shows that retrieval is the decisive component (+0.20 fact recall) and that BM25 beats dense
retrieval on this entity-heavy corpus. Multi-agent decomposition finds 17\% more of the referenced articles.
However, a single strong ReAct agent writes equally good answers at one eighth of the cost, and STE answers are far
more readable but less complete. More agents are not automatically better: each component has to earn its cost on
the target task, and only careful ablation shows which ones do.

\section{Reflection}
%%REFLECTION%%

\section*{Acknowledgment}
Claude (Anthropic) was used as a coding assistant. All design decisions, experiments and conclusions are my own and
were checked against the notebook outputs.
